
Par
passage aux gradués associés dans (3.2.1) et (3.2.2) on obtient un homomorphisme de degré \(-r\) :
\begin{equation}\tag{3.2.3}
 f_*:\Gro(X')_{\mathbb Q}\longrightarrow\Gro(X)_{\mathbb Q}
\end{equation}
et grâce à (1.3.2) une formule de projection graduée
\[
 f_*\bigl(f^*(x)\cdot y\bigr)=x\cdot f_*(y)
 \qquad (x\in\Gro(X)_{\mathbb Q},\ y\in\Gro(X')_{\mathbb Q}).
\]

\subsection*{3.3.}
Considérons enfin un morphisme \(f:X'\to X\) de Tor-dimension finie de schémas noethériens. Alors on sait (Exp. IV) que \(f\) induit un homomorphisme
\[
 f^*:K_\bullet(X)\longrightarrow K_\bullet(X')
\]
tel que
\[
 f^*(\cl_\bullet(\underline{F}))=
 \sum_i(-1)^i\cl_\bullet\bigl(\underline{\Tor}_i^{\underline{O}_X}(\underline{F},\underline{O}_{X'})\bigr)
\]
pour tout \(\underline{O}_X\)-module cohérent \(\underline{F}\).

Les \(\underline{O}_X\)-modules \(\underline{\Tor}_i^{\underline{O}_X}(\underline{F},\underline{O}_{X'})\) sont concentrés sur \(f^{-1}(\Supp(\underline{F}))\) et lorsque les fibres de \(f\) sont partout de dimension \(\leq d\), alors on sait (EGA IV 5.6.7) que
\[
 \dim(f^{-1}(\Supp(\underline{F})))\leq \dim(\Supp(\underline{F}))+d.
\]
Il s'ensuit que \(f\) induit pour tout \(k\in\mathbb Z\) un homomorphisme
\[
 f^*:\Filt_k(X)\longrightarrow\Filt_{k+d}(X'),
\]
d'où un homomorphisme de degré \(d\) de groupes abéliens gradués
